% GENERATED FILE. Edit canonical lessons, then run npm run generate:books.
% canonical-source-hash: fnv1a64:12440db94f563644
% canonical-lessons: SA-C424-paryatati, SA-C424-darsaniyah, SA-C424-ramaniyah, SA-C424-drsyam, SA-C424-viharati

\chapter{Sights Worth Seeing}
\label{ch:sa-sights-worth-seeing}
\hlchaptermodality{424}

\begin{chapteropening}
\textbf{By the end of this chapter:} \emph{I can say which places are worth seeing and why.}

Complete the last lesson of chapter 424: I can say which places are worth seeing and why.
\end{chapteropening}

\section[\texorpdfstring{\sk{पर्यटति} (paryaṭati)}{पर्यटति (paryaṭati)}]{\sk{पर्यटति} (paryaṭati) — tours; travels around}

\label{lesson:SA-C424-paryatati}



\begin{quote}
Before the new one: say the Sanskrit for a gap, then the Sanskrit for a gatekeeper.
\end{quote}

\subsection*{You'll want to know: \sk{पर्यटति}}
\textbf{\sk{पर्यटति}} — \emph{paryaṭati} — \textquotedblleft{}tours; travels around\textquotedblright{}.

\begin{quote}
\sk{पर्यटकः} \sk{नगरे} \sk{पर्यटति।}
\end{quote}

\begin{quote}
\emph{paryaṭakaḥ nagare paryaṭati.}
\end{quote}

\begin{quote}
The tourist tours the city.
\end{quote}

Past (he …): \sk{सः} \sk{नगरे} \sk{पर्यटितवान्}, \emph{saḥ nagare paryaṭitavān}.

\begin{cousinweb}
\textbf{\sk{पर्यटति}} (\emph{paryaṭati}) is built from \emph{pari}, \textquotedblleft{}around\textquotedblright{}, and the root √aṭ, \textquotedblleft{}to roam\textquotedblright{}, which also gives \emph{paryaṭaka}, \textquotedblleft{}a tourist\textquotedblright{}; MW gives \emph{pary-aṭ} as to roam about, travel over, with the participle \emph{paryaṭita} (MW).
\end{cousinweb}

\begin{grammarlens}[title={What you've built}]
One more word for getting around.
\end{grammarlens}

\subsection*{Your turn}
\begin{itemize}
\raggedright
  \item \emph{Say it:} \emph{paryaṭati}
  \item \emph{Say it:} \emph{paryaṭati}, once more
  \item \emph{Say it:} \emph{paryaṭakaḥ nagare paryaṭati.} — The tourist tours the city.
  \item \emph{Recall:} say the Sanskrit for a gap, then the Sanskrit for a gatekeeper, then say \emph{paryaṭati} again
\end{itemize}

\begin{tcolorbox}[breakable,colback=teal!4,colframe=teal!35!black,title={Before you move on}]
What does \sk{पर्यटति} mean? (\textquotedblleft{}Tours; travels around\textquotedblright{}.) Say it once more.
\end{tcolorbox}
\section[\texorpdfstring{\sk{दर्शनीयः} (darśanīyaḥ)}{दर्शनीयः (darśanīyaḥ)}]{\sk{दर्शनीयः} (darśanīyaḥ) — worth seeing}

\label{lesson:SA-C424-darsaniyah}



\begin{quote}
Before the new one: say the Sanskrit for a gatekeeper, then the Sanskrit for tours.
\end{quote}

\subsection*{You'll want to know: \sk{दर्शनीयः}}
\textbf{\sk{दर्शनीयः}} — \emph{darśanīyaḥ} — \textquotedblleft{}worth seeing\textquotedblright{}.

\begin{quote}
\sk{नगरस्य} \sk{प्राचीनः} \sk{प्रासादः} \sk{दर्शनीयः} \sk{अस्ति।}
\end{quote}

\begin{quote}
\emph{nagarasya prācīnaḥ prāsādaḥ darśanīyaḥ asti.}
\end{quote}

\begin{quote}
The city's old palace is worth seeing.
\end{quote}

\begin{cousinweb}
\textbf{\sk{दर्शनीयः}} (\emph{darśanīyaḥ}) is built on the root √dṛś, \textquotedblleft{}to see\textquotedblright{}, with the ending \emph{-anīya}, \textquotedblleft{}to be done\textquotedblright{}: to be seen; MW gives \emph{darśanīya} as worthy of being seen, beautiful (MW).
\end{cousinweb}

\begin{grammarlens}[title={What you've built}]
One more word for getting around.
\end{grammarlens}

\subsection*{Your turn}
\begin{itemize}
\raggedright
  \item \emph{Say it:} \emph{darśanīyaḥ}
  \item \emph{Say it:} \emph{darśanīyaḥ}, once more
  \item \emph{Say it:} \emph{nagarasya prācīnaḥ prāsādaḥ darśanīyaḥ asti.} — The city's old palace is worth seeing.
  \item \emph{Recall:} say the Sanskrit for a gatekeeper, then the Sanskrit for tours, then say \emph{darśanīyaḥ} again
\end{itemize}

\begin{tcolorbox}[breakable,colback=teal!4,colframe=teal!35!black,title={Before you move on}]
What does \sk{दर्शनीयः} mean? (\textquotedblleft{}Worth seeing\textquotedblright{}.) Say it once more.
\end{tcolorbox}
\section[\texorpdfstring{\sk{रमणीयः} (ramaṇīyaḥ)}{रमणीयः (ramaṇīyaḥ)}]{\sk{रमणीयः} (ramaṇīyaḥ) — delightful; lovely}

\label{lesson:SA-C424-ramaniyah}



\begin{quote}
Before the new one: say the Sanskrit for tours, then the Sanskrit for worth seeing.
\end{quote}

\subsection*{You'll want to know: \sk{रमणीयः}}
\textbf{\sk{रमणीयः}} — \emph{ramaṇīyaḥ} — \textquotedblleft{}delightful; lovely\textquotedblright{}.

\begin{quote}
\sk{सरोवरस्य} \sk{तीरं} \sk{रमणीयम्} \sk{अस्ति।}
\end{quote}

\begin{quote}
\emph{sarovarasya tīraṃ ramaṇīyam asti.}
\end{quote}

\begin{quote}
The lake shore is delightful.
\end{quote}

\begin{cousinweb}
\textbf{\sk{रमणीयः}} (\emph{ramaṇīyaḥ}) is built on the root √ram, \textquotedblleft{}to delight\textquotedblright{}, which also gives \emph{ramate}, \textquotedblleft{}enjoys\textquotedblright{}, with the ending \emph{-anīya}; MW gives \emph{ramaṇīya} as to be enjoyed, pleasant, delightful (MW).
\end{cousinweb}

\begin{grammarlens}[title={What you've built}]
One more word for getting around.
\end{grammarlens}

\subsection*{Your turn}
\begin{itemize}
\raggedright
  \item \emph{Say it:} \emph{ramaṇīyaḥ}
  \item \emph{Say it:} \emph{ramaṇīyaḥ}, once more
  \item \emph{Say it:} \emph{sarovarasya tīraṃ ramaṇīyam asti.} — The lake shore is delightful.
  \item \emph{Recall:} say the Sanskrit for tours, then the Sanskrit for worth seeing, then say \emph{ramaṇīyaḥ} again
\end{itemize}

\begin{tcolorbox}[breakable,colback=teal!4,colframe=teal!35!black,title={Before you move on}]
What does \sk{रमणीयः} mean? (\textquotedblleft{}Delightful; lovely\textquotedblright{}.) Say it once more.
\end{tcolorbox}
\section[\texorpdfstring{\sk{दृश्यम्} (dṛśyam)}{दृश्यम् (dṛśyam)}]{\sk{दृश्यम्} (dṛśyam) — a view; a sight}

\label{lesson:SA-C424-drsyam}



\begin{quote}
Before the new one: say the Sanskrit for worth seeing, then the Sanskrit for delightful.
\end{quote}

\subsection*{You'll want to know: \sk{दृश्यम्}}
\textbf{\sk{दृश्यम्}} — \emph{dṛśyam} — \textquotedblleft{}a view; a sight\textquotedblright{}.

\begin{quote}
\sk{पर्वतात्} \sk{नगरस्य} \sk{दृश्यं} \sk{सुन्दरम्} \sk{अस्ति।}
\end{quote}

\begin{quote}
\emph{parvatāt nagarasya dṛśyaṃ sundaram asti.}
\end{quote}

\begin{quote}
The view of the city from the mountain is beautiful.
\end{quote}

A neuter noun.

\begin{cousinweb}
\textbf{\sk{दृश्यम्}} (\emph{dṛśyam}) is built on the root √dṛś, \textquotedblleft{}to see\textquotedblright{}, with the ending \emph{-ya}: MW gives \emph{dṛśya} as visible, worth seeing and, as a noun, any visible object; a view is that sense (MW).
\end{cousinweb}

\begin{grammarlens}[title={What you've built}]
One more word for getting around.
\end{grammarlens}

\subsection*{Your turn}
\begin{itemize}
\raggedright
  \item \emph{Say it:} \emph{dṛśyam}
  \item \emph{Say it:} \emph{dṛśyam}, once more
  \item \emph{Say it:} \emph{parvatāt nagarasya dṛśyaṃ sundaram asti.} — The view of the city from the mountain is beautiful.
  \item \emph{Recall:} say the Sanskrit for worth seeing, then the Sanskrit for delightful, then say \emph{dṛśyam} again
\end{itemize}

\begin{tcolorbox}[breakable,colback=teal!4,colframe=teal!35!black,title={Before you move on}]
What does \sk{दृश्यम्} mean? (\textquotedblleft{}A view; a sight\textquotedblright{}.) Say it once more.
\end{tcolorbox}
\section[\texorpdfstring{\sk{विहरति} (viharati)}{विहरति (viharati)}]{\sk{विहरति} (viharati) — strolls; walks for pleasure}

\label{lesson:SA-C424-viharati}



\begin{quote}
Before the new one: say the Sanskrit for delightful, then the Sanskrit for a view.
\end{quote}

\subsection*{You'll want to know: \sk{विहरति}}
\textbf{\sk{विहरति}} — \emph{viharati} — \textquotedblleft{}strolls; walks for pleasure\textquotedblright{}.

\begin{quote}
\sk{सायं} \sk{वयम्} \sk{उद्याने} \sk{विहरामः।}
\end{quote}

\begin{quote}
\emph{sāyaṃ vayam udyāne viharāmaḥ.}
\end{quote}

\begin{quote}
In the evening we stroll in the garden.
\end{quote}

Past (he …): \sk{सः} \sk{उद्याने} \sk{विहृतवान्}, \emph{saḥ udyāne vihṛtavān}.

\begin{cousinweb}
\textbf{\sk{विहरति}} (\emph{viharati}) is built from \emph{vi}, \textquotedblleft{}apart\textquotedblright{}, and the root √hṛ, \textquotedblleft{}to carry\textquotedblright{}, which also gives \emph{harati}, \textquotedblleft{}takes away\textquotedblright{}; MW gives \emph{vi-hṛ} as to walk for pleasure, to roam, with the participle \emph{vihṛta} (MW).
\end{cousinweb}

\begin{grammarlens}[title={What you've built}]
One more word for getting around.
\end{grammarlens}

\subsection*{Your turn}
\begin{itemize}
\raggedright
  \item \emph{Say it:} \emph{viharati}
  \item \emph{Say it:} \emph{viharati}, once more
  \item \emph{Say it:} \emph{sāyaṃ vayam udyāne viharāmaḥ.} — In the evening we stroll in the garden.
  \item \emph{Recall:} say the Sanskrit for delightful, then the Sanskrit for a view, then say \emph{viharati} again
\end{itemize}

\begin{tcolorbox}[breakable,colback=teal!4,colframe=teal!35!black,title={Before you move on}]
What does \sk{विहरति} mean? (\textquotedblleft{}Strolls; walks for pleasure\textquotedblright{}.) Say it once more.
\end{tcolorbox}
